%%% mode: latex
%%% TeX-master: t
%%% End:

\chapter{基于通用基准的单体模型对抗鲁棒性评估}
\label{cha:single_agent}

\section{实验设置}
\label{sec:exp_setup}

\subsection{数据集}
\label{subsec:datasets}

为全面评估医疗大语言模型在不同攻击场景下的安全性表现，本研究在数据集选取上遵循以下原则：其一，功能覆盖的完备性，所选数据集须覆盖第二章所分析的数据污染与提示词操纵两大类共四种攻击方式，并为模型正常性能提供参照基准；其二，领域代表性，数据集应来源于权威机构或经同行评审的公开资源，能够反映真实医疗场景中的知识分布与任务特征；其三，评估维度的互补性，不同数据集之间应在问题类型、学科范围及任务形式上形成差异化覆盖，以避免单一基准带来的评估偏差。

基于上述原则，本研究选取的数据集按实验功能划分为三类。第一类为攻击载体数据集，包括用于模型微调训练及数据投毒实验的MedInstruct-52k\cite{zhang2023alpacare}，以及用于后门攻击效果测试的MedQuAD\cite{ben2019medquad}，二者共同支撑数据污染类攻击实验的实施与评估。第二类为对抗性测试数据集，包括HarmfulQA\cite{bhardwaj2023red}与ToxicQAFinal\cite{toxicqafinal2024}，分别针对越狱攻击和提示词注入攻击提供恶意查询样本，用于检验模型在提示词操纵类攻击下的防御能力。第三类为性能基线数据集，涵盖iCliniq\cite{li2023chatdoctor}、MEDQA\cite{jin2021medqa}、HeadQA\cite{vilares2019headqa}、PubmedQA\cite{jin2019pubmedqa}及MEDMCQA\cite{pal2022medmcqa}五个广泛使用的医疗问答基准，覆盖在线医疗问诊、执业医师资格考试、医学文献理解等多种任务形式，用以建立模型在未受攻击状态下的正常性能参照，从而有效区分模型能力下降究竟源于对抗性攻击还是自身的知识局限。各数据集的样本规模、来源及具体用途等详细信息如表\ref{tab:datasets}所示。

\begin{table}[htbp]
\centering
\caption{实验数据集统计信息}
\label{tab:datasets}
\begin{tabular}{lllll}
\toprule
数据集名称 & 样本量 & 问题类型 & 来源 & 用途 \\
\midrule
MedQuAD\cite{ben2019medquad} & 47,457对 & 多科室医疗问答 & NIH & 后门攻击测试 \\
MedInstruct-52k\cite{zhang2023alpacare} & 52,000条 & 医疗指令-响应对 & AlpaCare & 投毒攻击训练 \\
HarmfulQA\cite{bhardwaj2023red} & 1,960条 & 有害查询 & Declare-lab & 越狱攻击测试 \\
ToxicQAFinal\cite{toxicqafinal2024} & 1,514条 & 有毒查询 & HuggingFace & 提示词注入测试 \\
iCliniq\cite{li2023chatdoctor} & 31,062对 & 在线医疗问诊 & iCliniq.com & 性能基线评估 \\
MEDQA\cite{jin2021medqa} & 12,723条 & 医学考试选择题 & USMLE & 性能基线评估 \\
HeadQA\cite{vilares2019headqa} & 2,742条 & 健康科学考试题 & 西班牙医学考试 & 性能基线评估 \\
PubmedQA\cite{jin2019pubmedqa} & 1,000对 & 生物医学文献问答 & PubMed & 性能基线评估 \\
MEDMCQA\cite{pal2022medmcqa} & 193,155条 & 医学选择题 & AIIMS \& NEET PG & 性能基线评估 \\
\bottomrule
\end{tabular}
\end{table}

\subsection{待测模型}
\label{subsec:models}

为系统性评估不同模型架构与训练策略下的安全性能，本研究在待测模型选取上遵循以下原则：其一，架构多样性，所选模型应覆盖当前主流的预训练架构，以分析架构差异对安全性的潜在影响；其二，领域代表性，模型应涵盖医疗专用模型与通用模型两大类别，以揭示领域特定化对模型安全性的影响；其三，规模可比性，所选模型的参数规模应处于可比范围内，以控制模型容量对实验结果的干扰。基于上述原则，本研究选取待测模型，按领域属性分为医疗专用模型与通用基线模型两类，各模型的详细规格信息如表\ref{tab:models}所示。

\begin{table}[htbp]
\centering
\caption{待测模型规格信息}
\label{tab:models}
\resizebox{\textwidth}{!}{
\begin{tabular}{llllll}
\toprule
模型名称 & 参数量 & 基座架构 & 训练方式 & 领域属性 & 实验角色 \\
\midrule
BioGPT\cite{luo2022biogpt} & 1.5B & GPT-2 & 生物医学语料预训练 & 医疗专用 & 投毒/后门攻击 \\
AlpaCare\cite{zhang2023alpacare} & 7B & LLaMA & 医疗指令微调 & 医疗专用 & 全部攻击类型 \\
Asclepius\cite{kweon2023asclepius} & 7B & LLaMA & 合成临床笔记微调 & 医疗专用 & 提示词操纵/越狱攻击 \\
BioMistral\cite{labrak2024biomistral} & 7B & Mistral & 生物医学继续预训练 & 医疗专用 & 提示词操纵/越狱攻击 \\
BioLlama3\cite{gachon2024biollama} & 8B & LLaMA 3 & 医疗数据集微调 & 医疗专用 & 投毒/多智能体攻击 \\
Llama2\cite{touvron2023llama2} & 7B & LLaMA 2 & 通用RLHF对齐 & 通用基线 & 提示词操纵/越狱攻击 \\
Llama3\cite{grattafiori2024llama3} & 8B & LLaMA 3 & 通用RLHF对齐 & 通用基线 & 多智能体攻击 \\
\bottomrule
\end{tabular}
}
\end{table}

医疗专用模型方面，本研究选取了五个具有不同架构特征和训练策略的代表性模型。BioGPT\cite{luo2022biogpt}是微软研究院开发的生物医学预训练生成模型，基于GPT-2架构在大规模PubMed文献语料上进行预训练，参数量为1.5B，专为生物医学文本生成与挖掘任务设计，在关系抽取、问答和文档分类等任务中表现优异。AlpaCare\cite{zhang2023alpacare}基于LLaMA-7B架构，通过在包含52,000条医疗指令-响应对的MedInstruct数据集上进行指令微调构建而成，具有较强的医疗指令遵循能力，是本研究中覆盖攻击类型最广的待测模型。Asclepius\cite{kweon2023asclepius}同样基于LLaMA-7B架构，其特色在于使用GPT-3.5生成的合成临床笔记进行训练，面向临床文本理解和生成场景，在临床笔记摘要和病历分析等任务中展现出较好的性能。BioMistral\cite{labrak2024biomistral}基于Mistral-7B架构，在PubMed Central全文语料上进行生物医学领域的继续预训练，通过领域适配策略获得了扎实的医学知识基础，在多个医学问答基准上表现优异。BioLlama3\cite{gachon2024biollama}是基于最新的LLaMA 3-8B架构进行医疗领域微调的模型，结合了第三代LLaMA模型强大的语言理解能力与医疗领域的专业知识，主要用于数据投毒攻击和多智能体系统安全评估实验。

通用基线模型方面，本研究选取了Meta公司开发的Llama2\cite{touvron2023llama2}和Llama3\cite{grattafiori2024llama3}作为对照。Llama2-7B在大规模通用语料上进行预训练，并经过基于人类反馈的强化学习（RLHF）对齐，具有相对完善的安全机制。Llama3-8B作为其第三代演进版本，在训练数据规模、上下文窗口和推理能力方面均有显著提升，安全对齐机制进一步增强。通用模型的引入为评估医疗专用模型的安全性提供了对照基准——若医疗专用模型的抗攻击能力弱于同架构的通用模型，则表明领域特定化过程可能削弱了模型的安全性，这对于医疗大语言模型的安全设计具有重要的参考意义。

需要指出的是，本研究所选模型的参数量集中在7B至8B区间（BioGPT为1.5B），这一选择基于以下考量。其一，受医疗机构计算资源有限、数据隐私保护对本地化部署的刚性需求以及推理延迟约束等因素制约，7B量级的开源模型是当前医疗场景中最具部署可行性的方案\cite{thirunavukarasu2023large,clusmann2023future}，也是现有开源医疗大语言模型的主流参数规模，对该量级模型的安全性评估具有直接的现实意义。其二，本研究引入了同参数量的通用模型（Llama2、Llama3）作为对照基线，通过控制模型规模这一变量，可以有效区分安全性差异究竟源于模型容量不足还是领域适配过程本身的影响。其三，目前基于70B及以上规模构建的开源医疗专用模型数量有限，难以形成有效的对比实验组，更大规模模型的安全性评估将作为未来工作予以拓展。

\subsection{评价指标及其应用方式}
\label{subsec:metrics}

本章实验采用第二章第\ref{sec:eval_metrics}节所建立的评价指标体系，包括攻击成功率（ASR）、文本困惑度（PPL）和连贯性（COH）三个维度。各指标的定义及计算公式详见公式(\ref{eq:asr})、(\ref{eq:ppl})和(\ref{eq:coherence})，本节重点说明这些通用指标在本章不同攻击场景下的具体操作化方式。

ASR作为衡量模型对抗鲁棒性的核心指标，其判定标准需根据攻击类型做出针对性调整。在数据投毒攻击实验中，ASR定义为目标模型在受污染数据上微调后，对测试查询输出包含预设错误信息的比例，例如针对特定疾病查询时推荐攻击者预设的不当药物。在提示词注入攻击实验中，ASR定义为模型在接收含恶意注入指令的输入后，未能维持原始任务目标或遵从了恶意指令的比例。在越狱攻击实验中，ASR定义为模型未能对恶意查询生成适当拒绝响应的比例，即模型直接回答了应当被安全机制拦截的有害请求。上述三类实验中，ASR的判定均由MLSE框架中的安全智能体和评估智能体协同完成，基于关键词匹配与语义相似度计算的综合判定策略实现自动化评估。

PPL和COH在本章中承担攻击隐蔽性的诊断功能，用于回答一个关键问题：攻击是否在不引起文本质量异常的前提下成功操纵了模型输出。具体而言，对于每组攻击实验，本研究将攻击前后模型输出的PPL和COH进行对比分析。若攻击成功（ASR较高）但PPL和COH保持在模型正常运作范围内，则表明攻击具有高隐蔽性，此类攻击在医疗场景中尤为危险，因为流畅且连贯的错误输出更难被临床工作者识别。反之，若攻击导致PPL显著升高或COH明显下降，则表明攻击虽可能成功但易于通过文本质量监控手段检测。

三个指标之间形成互补的评估逻辑：ASR量化攻击的有效性，PPL反映输出的表面流畅度，COH衡量输入与输出之间的语义关联程度。三者结合能够全面刻画攻击的``有效性-隐蔽性''特征，为后续分析各类攻击对医疗大语言模型的实际威胁程度提供多维度的量化依据。

\subsection{实验环境与可复现性}
\label{subsec:exp_environment}

本研究所有实验均在同一软硬件环境下完成，具体配置如表\ref{tab:exp_env}所示。

\begin{table}[htbp]
\centering
\caption{实验软硬件配置}
\label{tab:exp_env}
\begin{tabular}{cc}
\toprule
配置项 & 具体参数 \\
\midrule
CPU & Intel(R) Xeon(R) Platinum 8268 (96 vCPUs) \\
GPU & $6 \times$ NVIDIA A100 SXM4 80GB \\
系统内存 & 1.0 TiB RAM \\
操作系统 & Ubuntu 24.04.3 LTS \\
Python & 3.10 \\
PyTorch & 2.0 \\
AutoGen & 0.2 \\
Ollama & 0.1 \\
\bottomrule
\end{tabular}
\end{table}

为保证实验的可复现性，所有实验统一使用固定随机种子42，模型推理温度设置为0。上述超参数在全部实验中保持一致，以排除随机性与生成策略差异对实验结果的干扰。

\section{数据投毒攻击实验分析}
\label{sec:poisoning}

\subsection{投毒比例对ASR的影响}
\label{subsec:poisoning_ratio}

数据投毒攻击通过在模型训练数据中注入精心设计的恶意样本，使目标模型在特定输入条件下产生攻击者预设的错误输出\cite{alber2024medical,shu2023exploit}。本研究采用AutoPoison工具\cite{shu2023exploit}自动生成投毒样本，系统性地考察投毒比例这一核心变量对攻击成功率的影响规律。实验选取BioGPT、AlpaCare和BioLlama3三个医疗大语言模型作为测试对象，以MedInstruct-52k为训练数据集，分别注入1\%、2\%、5\%和10\%四个比例的投毒数据进行微调，随后在MedQuAD测试集上评估各模型的ASR，具体结果如表\ref{tab:poison_results}所示。

\begin{table}[htbp]
\centering
\caption{数据投毒攻击实验结果}
\label{tab:poison_results}
\begin{tabular}{cccc}
\toprule
投毒比例 & BioGPT ASR(\%) & AlpaCare ASR(\%) & BioLlama3 ASR(\%) \\
\midrule
1\% & 0.00 & 1.24 & 0.00 \\
2\% & 0.00 & 1.48 & 0.00 \\
5\% & 0.00 & 8.44 & 0.00 \\
10\% & 0.13 & 12.86 & 0.14 \\
\bottomrule
\end{tabular}
\end{table}

实验结果表明，投毒比例与ASR之间存在明确的正相关关系，且不同模型对投毒攻击的敏感程度呈现显著差异。AlpaCare模型表现出最高的脆弱性，其ASR随投毒比例的增加呈现近似单调递增趋势：在1\%投毒比例下ASR为1.24\%，2\%时升至1.48\%，5\%时骤增至8.44\%，10\%时进一步攀升至12.86\%。相比之下，BioGPT和BioLlama3模型对投毒攻击的抵御能力显著优于AlpaCare，二者在1\%至5\%的投毒比例范围内ASR均保持为零，仅在投毒比例上升至10\%时方出现极低的ASR（分别为0.13\%和0.14\%）。上述差异初步表明，模型架构与训练策略可能影响其对数据投毒的敏感程度，但确切的影响机制尚需进一步的消融研究予以验证。

上述结果揭示了一个值得警惕的安全隐患：即便投毒数据仅占训练集的极小比例，亦足以导致特定模型在医疗问答任务中产生系统性错误。在临床应用场景中，此类错误可能直接转化为错误的诊断建议或不当的用药推荐，进而引发严重的医疗风险。值得注意的是，传统文本质量指标PPL和COH在投毒攻击前后并未出现显著波动，始终保持在模型的正常运作范围内。这一发现表明，常规的文本质量评估方法无法有效检测由数据投毒引入的医疗谬误，凸显了开发专门面向医疗场景的安全性评估手段的紧迫性。

为进一步验证数据投毒攻击的隐蔽性，即攻击是否会损害模型的通用基础能力，本研究引入MT-Bench基准对受攻击模型进行综合性能评估。考虑到评估成本与逻辑代表性，本节选取在投毒攻击中表现出最高脆弱性的AlpaCare模型作为代表性测试对象，对比不同投毒比例下该模型与未受攻击基线状态（Clear）的MT-Bench评分，结果如表\ref{tab:poison_mt_bench}所示。实验结果表明，各投毒比例下AlpaCare模型的GPT-4评分与GPT-3.5评分均未出现显著下降，投毒攻击并未导致模型通用能力的明显退化。这一现象进一步验证了投毒攻击的``静默失效''特征——此类攻击能够在不损害模型常规问答能力的前提下，精准地植入医学谬误，从而难以被传统的综合能力评估基准所察觉。

\begin{table}[htbp]
\centering
\caption{被投毒 AlpaCare 模型在 MT-Bench 上的性能评估}
\label{tab:poison_mt_bench}
\begin{tabular}{ccccccc}
\toprule
\multirow{2}{*}{投毒比例} & \multicolumn{3}{c}{GPT-4 评分} & \multicolumn{3}{c}{GPT-3.5 评分} \\
\cmidrule(lr){2-4} \cmidrule(lr){5-7}
 & 第一轮 & 第二轮 & 平均 & 第一轮 & 第二轮 & 平均 \\
\midrule
Clear & 5.87 & 4.56 & 5.22 & 6.78 & 6.22 & 6.51 \\
1\% & 5.54 & 3.69 & 4.62 & 6.34 & 6.11 & 6.23 \\
2\% & 5.96 & 4.11 & 5.03 & 6.33 & 6.21 & 6.27 \\
5\% & 5.87 & 4.24 & 5.05 & 6.46 & 5.96 & 6.22 \\
10\% & 6.03 & 3.96 & 4.99 & 6.64 & 6.49 & 6.57 \\
\bottomrule
\end{tabular}
\end{table}


\subsection{触发词的敏感性分析}
\label{subsec:trigger_sensitivity}

后门攻击作为数据投毒攻击的一种特殊形式，其核心机制在于向模型中植入隐性触发条件：受攻击模型在常规输入下仍可正常运作，仅当输入包含预设触发词时方激活恶意行为，输出攻击者预先设定的错误响应。本研究采用基于同义词词典的触发词设计方法\cite{chen2020badnl}，选取``drug''（药物）、``method''（方法）和``treatment''（治疗）三个在医疗文本中高频出现的术语作为触发词。实验通过同义词替换策略将上述触发词注入MedInstruct-52k数据集的样本中以构建后门攻击样本，并将其与正常样本按5\%的后门比例混合后分别微调AlpaCare和BioGPT模型，最终在MedQuAD数据集上评估不同触发词的激活效果。

实验结果如表\ref{tab:trigger_asr}所示，两个模型对后门攻击的敏感程度呈现显著差异。AlpaCare模型表现出较高的脆弱性，其中``treatment''触发词的ASR最高，达到3.18\%，``method''次之，为2.77\%，``drug''最低，为1.28\%。``treatment''一词触发效果最为显著这一现象，可能与其在医疗语境中的使用频率较高有关——高频词汇更容易在微调过程中被模型内化为有效的触发信号。相比之下，BioGPT模型对后门攻击展现出显著更强的抵御能力，三个触发词对应的ASR均处于0.12\%至0.14\%的极低水平，表明该模型可能具备某种内在的防御特性。值得关注的是，即便在后门攻击成功触发的情况下，两个模型输出的PPL和COH指标仍保持在正常范围内，表明后门攻击并未对生成文本的表面质量产生可检测的影响。

\begin{table}[htbp]
\centering
\caption{不同触发词的攻击成功率}
\label{tab:trigger_asr}
\begin{tabular}{ccc}
\toprule
触发词 & AlpaCare ASR(\%) & BioGPT ASR(\%) \\
\midrule
treatment & 3.18 & 0.14 \\
drug & 1.28 &  0.14 \\
method & 2.77 & 0.12 \\
\bottomrule
\end{tabular}
\end{table}

为进一步验证后门攻击的隐蔽性，即攻击是否会损害模型的通用基础能力，本研究沿用前节数据投毒实验的验证思路，选取脆弱性最为显著的AlpaCare模型作为代表性测试对象，在MT-Bench基准上对比不同触发词条件下受攻击模型与未受攻击基线模型（Clear）的综合性能评分。如表\ref{tab:backdoor_mt_bench}所示，无论采用何种触发词，受攻击模型的GPT-4评分和GPT-3.5评分均未出现显著下降，与基线模型基本持平。这一结果表明，后门攻击能够在不损害模型常规问答能力的前提下精准植入医学谬误，进一步印证了此类攻击在正常工作环境下具有高度隐蔽性、在特定触发条件下具有高危攻击性的双重特征。这种特性使得传统的模型测试规范及通用评价基准难以在此类攻击进入临床部署前实现有效检测，对医疗大语言模型的安全审计提出了严峻挑战。受计算资源限制，本研究目前选取三个具有代表性的医学触发词进行初步验证，未来工作将扩展至更大规模的触发词矩阵评估，以探究触发词语义类别与攻击效果之间的深层关联。

\begin{table}[htbp]
\centering
\caption{后门 AlpaCare 模型在 MT-Bench 上的性能评估}
\label{tab:backdoor_mt_bench}
\begin{tabular}{ccccccc}
\toprule
\multirow{2}{*}{触发词} & \multicolumn{3}{c}{GPT-4 评分} & \multicolumn{3}{c}{GPT-3.5 评分} \\
\cmidrule(lr){2-4} \cmidrule(lr){5-7}
 & 第一轮 & 第二轮 & 平均 & 第一轮 & 第二轮 & 平均 \\
\midrule
Clear & 5.87 & 4.56 & 5.22 & 6.78 & 6.22 & 6.51 \\
treatment & 5.99 & 3.85 & 4.92 & 6.71 & 6.39 & 6.56 \\
drug & 6.00 & 3.80 & 4.90 & 6.49 & 6.19 & 6.35 \\
method & 6.08 & 4.38 & 5.23 & 6.56 & 6.17 & 6.38 \\
\bottomrule
\end{tabular}
\end{table}

% 图4-2占位符：触发词敏感性对比柱状图
% TODO: 绘制柱状图对比不同触发词的攻击效果

\subsection{触发词长度与位置影响}
\label{subsec:trigger_length_position}

前节实验验证了不同触发词对后门攻击效果的影响，本节进一步从触发词长度和触发词插入位置两个维度，系统性地考察后门攻击的敏感性特征。

在触发词长度实验中，本研究设计了词长分别为1、2、3和4的触发词序列。针对AlpaCare模型，各长度对应的触发词依次为``method''、``therapeutic method''、``therapeutic method needed''和``therapeutic method we need''；针对BioGPT模型，各长度对应的触发词依次为``treatment''、``therapeutic treatment''、``therapeutic treatment needed''和``therapeutic treatment we need''。其他实验条件与前节保持一致，后门比例设定为5\%，实验结果如表\ref{tab:trigger_length}所示。对于AlpaCare模型，单词触发词（L=1）的ASR最高，达到3.18\%，双词触发词（L=2）骤降至1.28\%，三词触发词（L=3）回升至2.77\%，四词触发词（L=4）再次下降至1.38\%，整体呈现随长度增加而下降的趋势。这一现象表明，较短的触发词在微调过程中更易被模型作为统一的语义单元内化，从而建立起更稳定的触发-响应映射关系；而较长的触发词序列则因词间组合方式的多样性，增加了模型精确识别完整触发模式的难度，致使后门激活概率降低。相比之下，BioGPT模型在四种长度条件下的ASR均维持在0.12\%至0.17\%的极低水平，未呈现出与触发词长度相关的规律性变化，进一步印证了该模型对后门攻击所具备的较强内在鲁棒性。

\begin{table}[htbp]
\centering
\caption{触发词长度对攻击效果的影响}
\label{tab:trigger_length}
\begin{tabular}{ccc}
\toprule
模型 & 触发词长度 & ASR(\%) \\
\midrule
\multirow{4}{*}{AlpaCare} & L=1 & 3.18 \\
 & L=2 & 1.28 \\
 & L=3 & 2.77 \\
 & L=4 & 1.38 \\
\midrule
\multirow{4}{*}{BioGPT} & L=1 & 0.14 \\
 & L=2 & 0.14 \\
 & L=3 & 0.12 \\
 & L=4 & 0.17 \\
\bottomrule
\end{tabular}
\end{table}

在触发词位置实验中，本研究分别将触发词置于输入提示词的起始位置（Start）、中间位置（Middle）和末尾位置（End）进行测试，实验结果如表\ref{tab:trigger_position}所示。对于AlpaCare模型，末尾位置的ASR最高，为3.18\%，起始位置次之，为3.12\%，中间位置最低，仅为2.53\%。末尾位置攻击效果的优势可归因于 Transformer 架构中自回归解码机制的特性\cite{chen2020badnl}：位于输入序列末端的词元与模型生成阶段在时序上最为接近，在解码过程中对输出分布的影响更为直接，因而更易触发预植入的后门行为。中间位置的触发词则容易被前后上下文信息所稀释，其触发信号在注意力分配中的显著性相应降低。BioGPT模型在三种位置条件下的ASR均稳定在0.12\%至0.14\%之间，进一步验证了该模型对后门攻击的一致性抵御能力。

\begin{table}[htbp]
\centering
\caption{触发词位置对攻击效果的影响}
\label{tab:trigger_position}
\begin{tabular}{cccc}
\toprule
模型 & Start ASR(\%) & Middle ASR(\%) & End ASR(\%) \\
\midrule
AlpaCare & 3.12 & 2.53 & 3.18 \\
BioGPT & 0.14 & 0.12 & 0.12 \\
\bottomrule
\end{tabular}
\end{table}

综合触发词长度与位置两方面的实验结果，可以得出以下结论：在医疗场景的后门攻击中，短触发词与末尾插入位置的组合构成最具威胁的攻击配置，此类配置下后门激活概率最高且攻击隐蔽性最强。这一发现为医疗大语言模型的安全防御提供了明确的关注方向——在训练数据审计和输入过滤机制的设计中，应对输入序列尾部的短词汇异常模式予以重点监控。

\subsection{Oracle模型对投毒效果的影响}
\label{subsec:oracle_model}

AutoPoison工具依赖Oracle模型生成对抗性上下文和恶意响应，因此Oracle模型的能力直接影响投毒攻击的有效性。本研究对比了GPT-3.5和Gemma-7B两个不同规模的Oracle模型在投毒攻击中的表现差异。实验在AlpaCare和BioGPT两个目标模型上分别使用1\%、2\%、5\%和10\%的投毒比例进行测试，系统性地评估Oracle模型能力与攻击效果之间的关系。实验结果显示，使用GPT-3.5作为Oracle模型时，各投毒比例下的ASR均显著高于使用Gemma-7B的情况。以AlpaCare模型为例，当投毒比例为10\%时，使用GPT-3.5生成的投毒数据可使ASR达到12.86\%，而使用Gemma-7B时的ASR仅为3.34\%。这种显著差异表明Oracle模型的能力是影响投毒攻击效果的关键因素之一。

深入分析发现，GPT-3.5生成的投毒样本在质量和可读性方面明显优于Gemma-7B。Gemma-7B生成的投毒数据往往包含大量重复性句子，并且对不同问题生成相同或高度相似的答案，这种单一性显著降低了样本的质量和欺骗性。相比之下，GPT-3.5能够生成语义连贯、上下文相关的投毒样本，这些样本更接近真实医疗场景中的表述方式，因此更容易被目标模型学习并内化。为确保实验结果的可比性，除Oracle模型不同外，其他实验条件保持一致：投毒比例相同、目标模型相同、训练轮数相同（1 epoch）、随机种子固定为42。需要指出的是，GPT-3.5与Gemma-7B在模型规模、训练数据等方面存在差异，这些因素可能共同影响投毒效果，后续研究可通过更多Oracle模型对比来分离各因素影响。

% 图4-4占位符：不同Oracle模型投毒效果对比柱状图
% TODO: 绘制柱状图对比GPT-3.5与Gemma-7B的投毒效果

此外，本研究还探索了通过角色提示工程增强Oracle模型攻击效果的方法。具体而言，通过为Oracle模型分配``药剂师''角色，使其从医疗专业人员的角度生成投毒响应。实验结果表明，AutoPoison-Role方法的ASR达到15.77\%，明显高于基础AutoPoison方法的12.86\%。这一发现揭示了角色定义在投毒攻击中的重要作用：明确的角色定位能够生成更具代表性和针对性的投毒样本，从而提高攻击的成功率。具体 ASR 对比如表\ref{tab:poison_role_asr}所示。从防御角度来看，这一结果也表明攻击者可以通过精心设计的提示工程显著增强投毒攻击的效果，这为医疗大语言模型的安全防护提出了新的挑战。

\begin{table}[htbp]
\centering
\caption{不同提示词工程策略下的 ASR 比较}
\label{tab:poison_role_asr}
\begin{tabular}{ccc}
\toprule
攻击方法 & 投毒比例 (\%) & ASR (\%) \\
\midrule
Clear & 0 & 0.00 \\
AutoPoison & 10 & 12.86 \\
AutoPoison-Role & 10 & 15.77 \\
\bottomrule
\end{tabular}
\end{table}

此外，不同 Oracle 模型生成投毒数据的效果对比如表\ref{tab:poison_oracle_asr}所示。

\begin{table}[htbp]
\centering
\caption{不同 Oracle 模型对投毒效果的影响}
\label{tab:poison_oracle_asr}
\begin{tabular}{ccccc}
\toprule
\multirow{2}{*}{投毒比例} & \multicolumn{2}{c}{AlpaCare ASR(\%)} & \multicolumn{2}{c}{BioGPT ASR(\%)} \\
\cmidrule(lr){2-3} \cmidrule(lr){4-5}
 & GPT-3.5 & Gemma-7B & GPT-3.5 & Gemma-7B \\
\midrule
1\% & 1.24 & 1.38 & 0.00 & 0.00 \\
2\% & 1.48 & 1.26 & 0.00 & 0.00 \\
5\% & 8.44 & 1.32 & 0.00 & 0.00 \\
10\% & 12.86 & 3.34 & 0.13 & 0.10 \\
\bottomrule
\end{tabular}
\end{table}

\section{提示词操纵攻击实验分析}
\label{sec:prompt_injection}

\subsection{不同越狱策略的有效性对比}
\label{subsec:jailbreak_strategies}

提示词操纵攻击通过精心设计的恶意输入绕过模型的安全约束，迫使模型产生违背其原始训练目标的输出。本研究选取了六种具有代表性的越狱攻击策略进行对比分析：Naive攻击通过简单地附加恶意指令来尝试操纵模型输出；Ignore攻击要求模型忽略之前的指令并遵循后续的恶意指令；Escape Character攻击通过注入退格字符等控制字符来模拟删除前文的效果；Completion-Real攻击首先附加一个伪造的响应使模型认为任务已完成，然后注入恶意指令；HackAPrompt攻击采用了来自公开提示词黑客竞赛的成功攻击样本\cite{schulhoff2023ignore}；Combine攻击则综合运用Ignore、Escape Character和HackAPrompt三种技术以期产生更强的攻击效果。

实验在HarmfulQA和ToxicQAFinal两个对抗性数据集上对AlpaCare、Asclepius、BioMistral、BioLlama3、Llama2和Llama3模型进行测试。结果显示，医疗专用模型对提示词操纵攻击的脆弱性普遍高于通用Llama2和Llama3模型。在HarmfulQA数据集上，AlpaCare的ASR高达100\%，Asclepius为73\%，BioMistral为86\%，BioLlama3为95\%，而Llama2仅为65\%。具体结果框架如表\ref{tab:jailbreak_strategies}所示，详细数据待后续补充。

\begin{table}[htbp]
\centering
\caption{六种越狱策略攻击效果对比(ASR)}
\label{tab:jailbreak_strategies}
\begin{tabular}{lcccccc}
\toprule
攻击策略 & AlpaCare & Asclepius & BioMistral & BioLlama3 & Llama2 & Llama3 \\
\midrule
Naive & - & 72\% & - & - & - & - \\
Ignore & - & - & - & - & - & - \\
Escape Character & 100\% & - & - & - & - & - \\
Completion-Real & - & - & - & - & - & - \\
HackAPrompt & - & - & - & - & - & - \\
Combine & 100\% & - & - & - & - & - \\
\bottomrule
\end{tabular}
\end{table}
% TODO: 本表数据主要来自原图Fig 3c和3d，当前仅填入了正文中明确提及的数据，其余数据需从原始图表或实验记录中补充完整。

AlpaCare模型达到100\%的ASR这一异常结果值得深入分析。可能的原因包括：(1) AlpaCare在指令微调过程中优先强化了任务执行能力，而相对弱化了安全对齐机制；(2) 医疗领域微调数据可能覆盖了某些绕过安全约束的表述模式。这一发现揭示了医疗模型在专业化过程中可能存在的``安全性-专业性权衡''问题，为医疗大语言模型的安全设计提供了重要警示。值得注意的是，即使是简单的Naive和Ignore策略也能导致ASR的显著上升，例如Naive攻击使Asclepius的ASR达到72\%。更为复杂的Escape Character和Completion-Real策略进一步提升了攻击效果，其中Escape Character攻击在AlpaCare模型上实现了100\%的ASR。Combine攻击策略最为有效，在多个模型上均达到了接近或超过90\%的ASR。这些结果表明，医疗大语言模型在提示词操纵攻击面前表现出令人担忧的脆弱性，其安全对齐机制亟待加强。

\subsection{文本质量与攻击隐蔽性分析}
\label{subsec:ppl_stealth}

攻击的隐蔽性是评估对抗性攻击威胁程度的关键维度。一个具有高隐蔽性的攻击能够在不被察觉的情况下成功操纵模型输出，这种``静默失效''在医疗场景中尤其危险。困惑度作为衡量文本流畅度和自然性的指标，可以用来评估攻击的隐蔽性水平。本研究对不同提示词操纵攻击策略下模型输出的PPL进行了系统分析。实验结果显示，在大多数攻击策略下，模型输出的PPL并未出现显著波动，始终保持在模型的正常运作范围内。具体而言，Llama2模型在Escape Character和Combine攻击方法上表现出相对较高的PPL值，而这两种方法的ASR恰好较高，这表明Llama2可能对某些特定的攻击模式更敏感，从而触发了其安全机制产生更不自然的输出。

然而，对于其他大多数攻击方法，PPL结果并未呈现出明确的变化规律。这一发现揭示了对抗性攻击的一个令人担忧的特征：高攻击成功率并不一定伴随着文本质量的下降。换句话说，模型可以在生成流畅、语法正确、语义连贯的医疗谬误的同时，保持极低的PPL值。这种``自信的幻觉''实际上是最危险的情况，因为用户很难从输出文本本身判断其正确性，从而更容易被误导并采纳错误的医疗建议。从防御角度来看，传统的基于文本质量评估的检测方法在面对此类高隐蔽性攻击时效果有限，因此需要开发专门的医疗谬误检测机制，结合医学知识图谱和临床规则引擎，从语义层面而非表面特征层面识别模型输出的潜在风险。

% 图4-5占位符：攻击前后PPL分布箱线图
% TODO: 绘制箱线图对比不同攻击策略下的PPL分布

\subsection{提示词长度对攻击效果的影响}
\label{subsec:prompt_length}

提示词长度是影响攻击效果的重要因素之一。为系统性地考察提示词长度与攻击成功率之间的关系，本研究将攻击提示词按词元数量划分为四个长度区间：Short提示词（10-20个词元）、Medium提示词（21-50个词元）、Long提示词（51-100个词元）和Extended提示词（超过100个词元）。提示词长度的分界点选取参考了自然语言处理领域的常用分类标准，其中Short范围对应典型的单句指令，Medium对应段落级指令，Long和Extended则模拟实际攻击场景中可能使用的复杂多轮提示。实验分别使用HackAPrompt、Ignore和Escape Character三种攻击方法，在AlpaCare、Asclepius、BioMistral、BioLlama3、Llama2和Llama3模型上进行测试。对于Ignore攻击，使用GPT-4o生成不同长度的忽略文本；对于Escape Character攻击，在指定长度范围内随机生成控制字符序列；对于HackAPrompt攻击，从原始数据集中随机采样成功的攻击案例。

% 图4-6占位符：提示词长度×策略×模型雷达图
% TODO: 绘制雷达图展示多维度的攻击效果对比

\begin{table}[htbp]
\centering
\caption{HackAPrompt攻击在不同提示词长度下的ASR}
\label{tab:prompt_length_hackaprompt}
\begin{tabular}{ccccc}
\toprule
模型 & Short & Medium & Long & Extended \\
\midrule
AlpaCare & 0.37 & 0.32 & 0.28 & 0.46 \\
Asclepius & 0.23 & 0.34 & 0.51 & 0.58 \\
BioMistral & 0.23 & 0.19 & 0.22 & 0.16 \\
BioLlama3 & - & - & - & - \\
Llama2 & 0.55 & 0.52 & 0.54 & 0.50 \\
Llama3 & - & - & - & - \\
\bottomrule
\end{tabular}
\end{table}

\begin{table}[htbp]
\centering
\caption{Ignore攻击在不同提示词长度下的ASR}
\label{tab:prompt_length_ignore}
\begin{tabular}{ccccc}
\toprule
模型 & Short & Medium & Long & Extended \\
\midrule
AlpaCare & 0.18 & 0.13 & 0.06 & 0.10 \\
Asclepius & 0.35 & 0.40 & 0.18 & 0.22 \\
BioMistral & 0.12 & 0.22 & 0.18 & 0.09 \\
BioLlama3 & - & - & - & - \\
Llama2 & 0.81 & 0.70 & 0.21 & 0.04 \\
Llama3 & - & - & - & - \\
\bottomrule
\end{tabular}
\end{table}

\begin{table}[htbp]
\centering
\caption{Escape Character攻击在不同提示词长度下的ASR}
\label{tab:prompt_length_escape}
\begin{tabular}{ccccc}
\toprule
模型 & Short & Medium & Long & Extended \\
\midrule
AlpaCare & 0.45 & 0.26 & 0.18 & 0.25 \\
Asclepius & 0.30 & 0.33 & 0.39 & 0.09 \\
BioMistral & 0.11 & 0.07 & 0.03 & 0.02 \\
BioLlama3 & - & - & - & - \\
Llama2 & 0.81 & 0.70 & 0.21 & 0.04 \\
Llama3 & - & - & - & - \\
\bottomrule
\end{tabular}
\end{table}

实验结果表明，提示词长度对攻击效果的影响因攻击方法和模型而异。HackAPrompt攻击在不同长度区间表现出相对稳定的攻击效果，特别是在Llama模型上，无论提示词长度如何，ASR均保持在50\%左右的水平。Ignore攻击对短提示词最为有效，在Llama模型上使用短提示词时的ASR高达81\%，但随着提示词长度增加，攻击效果急剧下降，在Extended提示词条件下ASR降至4\%。Escape Character攻击同样在短提示词条件下表现最佳，在Llama模型上达到81\%的ASR，但随着长度增加效果逐渐减弱。一个值得注意的例外是BioMistral模型，该模型在所有攻击方法和长度条件下均表现出较强的抵御能力，ASR始终低于0.11。从安全角度来看，这些发现表明短提示词对大多数攻击类型而言风险最高，而不同模型对不同长度攻击的脆弱性差异显著。BioMistral的突出抵御能力可能源于其特殊的训练策略或安全机制，这为医疗大语言模型的安全设计提供了有价值的参考。

\subsection{提示词位置对攻击效果的影响}
\label{subsec:prompt_position}

恶意提示词在输入序列中的位置可能显著影响模型的响应模式。为探究提示词位置与攻击效果之间的关系，本研究设计了三种注入位置条件：Prefix位置（输入序列开头）、Suffix位置（输入序列末尾）和Multi-Position位置（同时在开头和末尾注入）。实验在AlpaCare、Asclepius、BioMistral、BioLlama3、Llama2和Llama3模型上进行，系统性地评估不同注入位置下的攻击成功率，具体结果如表\ref{tab:prompt_position}所示。

\begin{table}[htbp]
\centering
\caption{不同注入位置下的攻击成功率}
\label{tab:prompt_position}
\begin{tabular}{cccc}
\toprule
模型 & Prefix & Suffix & Multi \\
\midrule
AlpaCare & 0.2793 & 0.5602 & 0.4645 \\
Asclepius & 0.2894 & 0.5293 & 0.5247 \\
BioMistral & 0.4267 & 0.5682 & 0.4540 \\
BioLlama3 & - & - & - \\
Llama2 & 0.0934 & 0.1402 & 0.2202 \\
Llama3 & - & - & - \\
\bottomrule
\end{tabular}
\end{table}

实验结果显示，Suffix注入在所有测试模型上均表现出最高的攻击效果。BioMistral模型对Suffix注入最为脆弱，ASR达到0.5682，AlpaCare和Asclepius的ASR也分别达到0.5602和0.5293。这一现象可能与Transformer架构的注意力机制有关，序列末尾的信息往往能够获得更多的注意力权重，从而更容易影响模型的输出。相比之下，Prefix注入的攻击效果普遍较弱，BioMistral在Prefix条件下的ASR为0.4267，低于Suffix条件下的表现，但仍高于AlpaCare和Asclepius。Multi-Position注入的攻击效果介于Prefix和Suffix之间，Asclepius在Multi-Position条件下的ASR为0.5247。值得注意的是，Llama模型在所有注入位置条件下均表现出相对较强的抵御能力，ASR在0.0934至0.2202之间，显著低于医疗专用模型。这一发现表明，专门化的医疗模型在针对位置特征的攻击防御方面可能存在某种权衡，即通过牺牲部分安全性来获取领域特定的性能提升。从防御角度来看，加强模型对序列末尾输入的过滤和验证机制，是提升医疗大语言模型安全性的重要方向。

\subsection{攻击频率对成功率的影响}
\label{subsec:attack_frequency}

在实际攻击场景中，攻击者可能会在单次对话中多次尝试注入恶意提示词，以期增强攻击效果。本研究系统性地评估了注入频率对攻击成功率的影响，设计了1至7次重复注入的实验条件。对于每个目标模型和频率类别，我们在保持对话语境一致的前提下，嵌入相同内容的恶意提示词，并统计不同重复次数下的ASR变化，具体结果如表\ref{tab:attack_frequency}所示。

\begin{table}[htbp]
\centering
\caption{不同注入频率下的攻击成功率}
\label{tab:attack_frequency}
\begin{tabular}{cccccccc}
\toprule
模型 & 1次 & 2次 & 3次 & 4次 & 5次 & 6次 & 7次 \\
\midrule
AlpaCare & 0.4827 & 0.3025 & 0.2583 & 0.2164 & 0.1895 & 0.1787 & 0.1704 \\
Asclepius & 0.5682 & 0.5373 & 0.5182 & 0.4583 & 0.4120 & 0.3741 & 0.3926 \\
BioMistral & 0.5090 & 0.2380 & 0.1889 & 0.1852 & 0.1781 & 0.1630 & 0.1728 \\
BioLlama3 & - & - & - & - & - & - & - \\
Llama2 & 0.1630 & 0.1880 & 0.1926 & 0.1722 & 0.1355 & 0.1114 & 0.0914 \\
Llama3 & - & - & - & - & - & - & - \\
\bottomrule
\end{tabular}
\end{table}

实验结果显示，所有医疗专用模型在单次注入条件下表现出最高的脆弱性。Asclepius模型的初始ASR最高，达到0.5682，AlpaCare和BioMistral也分别达到0.4827和0.5090。随着注入频率增加，各模型的ASR整体呈下降趋势。BioMistral和AlpaCare在两次注入后ASR急剧下降，稳定在0.17左右的水平。Asclepius的下降趋势相对平缓，在七次注入后仍保持0.3926的ASR，表现出较强的持续性脆弱性。Llama模型呈现出独特的模式：其初始脆弱性较低（ASR为0.1630），随着注入频率增加，ASR略有上升后持续下降，始终保持着最低的整体脆弱性水平。

从机制角度分析，这种ASR随重复注入而下降的现象可能源于多个因素。首先是上下文窗口饱和效应，随着恶意提示词的重复出现，其在上下文中的相对重要性可能被稀释。其次是重复检测机制的触发，某些模型可能具备识别重复输入模式的能力，从而对重复内容降低响应权重。从防御角度来看，这些发现表明单次注入往往是最有效的攻击策略，重复注入的边际收益递减。因此，防御机制应当优先关注对初始注入的检测和拦截，同时建立基于频率的异常检测系统，以识别可能存在的攻击尝试。

\section{越狱攻击实验分析}
\label{sec:jailbreaking}

\subsection{基线模型直接恶意查询响应评估}
\label{subsec:baseline_malicious}

为评估模型在未采用特定攻击策略情况下对恶意查询的固有防御能力，本研究开展了直接恶意查询响应评估实验。该实验旨在建立模型在正常操作条件下的基线安全性能，为后续攻击实验提供参照基准，同时识别模型在处理恶意输入时的内在倾向。实验采用简化的评估标准：如果模型未能生成适当的拒绝响应（如``很抱歉，我无法满足您的请求''等），则视为攻击成功。这种评估方法能够揭示模型的基础防御机制及其在识别和拒绝潜在有害请求方面的有效性。

实验在HarmfulQA和ToxicQA两个数据集上对AlpaCare、Asclepius、BioMistral、BioLlama3、Llama2和Llama3模型进行测试，具体结果如表\ref{tab:baseline_response}所示。结果显示，不同模型对直接恶意查询的防御能力存在显著差异。AlpaCare模型表现出最高的脆弱性，在HarmfulQA和ToxicQA数据集上的ASR分别达到0.72和0.69，表明该模型的安全机制较为薄弱，倾向于直接响应有害查询。相比之下，Asclepius和BioMistral展现出强大的防御能力，在两个数据集上的ASR均为0.00，能够一致性地拒绝恶意查询。Llama模型的防御能力介于两者之间，在HarmfulQA和ToxicQA上的ASR分别为0.11和0.05。值得注意的是，Llama2模型采用了一种独特的教育性防御策略：面对恶意查询时，该模型并非直接拒绝，而是提供关于健康风险的教育性内容。这种设计虽然在一定程度上降低了直接危害，但也可能被用户理解为对查询的确认或部分认可。从医疗安全的角度来看，Asclepius和BioMistral的明确拒绝策略更为可靠，而AlpaCare的薄弱防御机制亟待加强。

\begin{table}[htbp]
\centering
\caption{模型对直接恶意查询的基线响应}
\label{tab:baseline_response}
\begin{tabular}{lccc}
\toprule
模型 & HarmfulQA ASR & ToxicQA ASR & 响应类型特征 \\
\midrule
AlpaCare & 0.72 & 0.69 & 直接回答 \\
Asclepius & 0.00 & 0.00 & 明确拒绝 \\
BioMistral & 0.00 & 0.00 & 明确拒绝 \\
BioLlama3 & - & - & - \\
Llama2 & 0.11 & 0.05 & 教育性回应 \\
Llama3 & - & - & - \\
\bottomrule
\end{tabular}
\end{table}

需要指出的是，实验采用简化的二元判定标准（拒绝/不拒绝）进行初步评估。这种标准可能遗漏部分边界情况，如模型提供``有条件的拒绝''或``教育性回应''。在第五章人机验证环节，我们将引入更细粒度的三级判定标准（明确拒绝/部分有害/完全有害），并由两位医学专家进行交叉验证，以提高评估的准确性。

\subsection{AutoDAN越狱攻击的跨样本泛化分析}
\label{subsec:autodan_generalization}

对抗性攻击的泛化能力是评估其现实威胁程度的关键指标。一个成功的攻击如果仅能在单个查询上生效，其危害相对有限；然而，如果攻击能够泛化到多个不同的查询，则可能揭示模型防御机制中的系统性漏洞。本研究基于AutoDAN（Automatic Distinctive and Natural Attack Generation）方法\cite{liu2023autodan}开展跨样本泛化分析，该方法采用分层遗传算法生成越狱提示词，通过段落级和句子级的结构变异来增强多样性并避免局部最优。实验设计为：对于每个成功的攻击提示词$i$，测试其在后续20个查询（$i+1$至$i+20$）上的泛化效果。选择20个后续样本作为泛化测试范围是基于以下考虑：(1) 在80\%置信水平下，该样本量足以检测15\%以上的泛化效应；(2) 该数量在计算成本与统计效力之间取得平衡；(3) 更大规模的泛化测试将作为未来工作拓展。

% 图4-7占位符：攻击效果随样本距离的变化趋势
% TODO: 绘制折线图展示泛化效果

实验结果如表\ref{tab:cross_sample}所示，不同模型对跨样本泛化攻击的抵御能力存在显著差异。AlpaCare模型表现出最高的泛化脆弱性，在HarmfulQA和ToxicQA数据集上的泛化ASR分别达到0.7540和0.6308，表明该模型的防御机制存在系统性弱点。Asclepius模型展现出最强的抵御能力，在两个数据集上的泛化ASR均低于0.06。BioMistral模型的抵御能力介于两者之间，对ToxicQA的泛化ASR（0.1193）明显高于对HarmfulQA的ASR（0.0715）。Llama模型表现出不对称的抵御模式：对HarmfulQA的泛化ASR为0.2592，但对ToxicQA完全免疫，泛化ASR为0.0000。

\begin{table}[htbp]
\centering
\caption{AutoDAN越狱攻击的跨样本泛化ASR}
\label{tab:cross_sample}
\begin{tabular}{ccc}
\toprule
模型 & HarmfulQA & ToxicQA \\
\midrule
AlpaCare & 0.7540 & 0.6308 \\
Asclepius & 0.0349 & 0.0523 \\
BioMistral & 0.0715 & 0.1193 \\
BioLlama3 & - & - \\
Llama2 & 0.2592 & 0.0000 \\
Llama3 & - & - \\
\bottomrule
\end{tabular}
\end{table}

从攻击类型的角度分析，HarmfulQA攻击的整体泛化效果优于ToxicQA，尤其是在AlpaCare模型上。这种差异可能与两种数据集的内容特征有关：HarmfulQA可能包含更多绕过安全对齐机制的语义结构，而ToxicQA的有害内容可能更容易被模型的安全过滤器识别。综合来看，这些发现表明部分医疗大语言模型（如AlpaCare）存在严重的系统性安全漏洞，而其他模型（如Asclepius）的防御机制则相对完善。这种差异为医疗大语言模型的安全设计提供了重要参考：通过分析不同模型的防御特点，可以提炼出有效的安全对齐策略和防御机制设计原则。

\subsection{重复攻击尝试次数分析}
\label{subsec:repeated_attempts}

在实际攻击场景中，攻击者可能通过多次尝试不同的越狱提示词来绕过模型的防御机制。为评估模型对重复攻击尝试的抵御能力，本研究设计了迭代测试协议：每个恶意查询使用提示词变异攻击进行最多100次尝试，当攻击成功或达到100次上限时结束实验。100次尝试上限设置参考了AutoDAN原论文的默认配置\cite{liu2023autodan}，该阈值足以覆盖99\%以上的成功案例（基于预实验统计）。实验统计了成功样本的平均尝试次数和所有样本的平均尝试次数两个指标，以全面评估模型的抵御特性。

% 图4-8占位符：成功攻击所需尝试次数分布直方图
% TODO: 绘制直方图展示不同模型的尝试次数分布

实验结果如表\ref{tab:repeated_attempts}所示，不同模型对重复攻击尝试的抵御能力存在显著差异。AlpaCare模型表现出最低的抵御能力，成功样本的平均尝试次数仅为1.5625次，且该数值与所有样本的平均尝试次数相同，表明该模型几乎在第一次尝试时就被成功攻破，防御机制近乎失效。Asclepius模型展现出最强的抵御能力，成功样本的平均尝试次数为28.9023次，当包含不成功样本时，平均尝试次数上升至47.21次，表明许多攻击尝试在经过多次尝试后仍然失败。BioMistral模型的抵御能力介于两者之间，成功样本的平均尝试次数为16.3532次，所有样本的平均尝试次数为23.045次。Llama模型表现出独特的抵御模式：成功样本的平均尝试次数相对较少（13.2727次），但所有样本的平均尝试次数高达71.38次，这表明虽然部分攻击能够相对快速地成功，但大量攻击尝试在经过数十次尝试后仍然失败。

\begin{table}[htbp]
\centering
\caption{重复攻击的平均尝试次数}
\label{tab:repeated_attempts}
\begin{tabular}{ccc}
\toprule
模型 & 成功样本平均次数 & 所有样本平均次数 \\
\midrule
AlpaCare & 1.5625 & 1.5625 \\
Asclepius & 28.9023 & 47.2100 \\
BioMistral & 16.3532 & 23.0450 \\
BioLlama3 & - & - \\
Llama2 & 13.2727 & 71.3800 \\
Llama3 & - & - \\
\bottomrule
\end{tabular}
\end{table}

从防御机制设计的角度来看，这些发现揭示了不同模型采用的防御策略差异。AlpaCare的薄弱防御可能源于其训练过程中对安全对齐的重视不足，或者其安全机制容易被精心设计的提示词绕过。Asclepius和BioMistral的较强抵御能力可能得益于更严格的安全对齐训练或专门的防御机制。Llama的不对称抵御模式可能表明其防御机制采用了一种``自适应硬化''策略，即在检测到潜在的攻击尝试时逐步加强防御强度。从攻击者的角度，这些结果表明针对不同模型需要采用不同的攻击策略：对于AlpaCare这类脆弱模型，简单的攻击提示词即可奏效；而对于Asclepius这类防御较强的模型，则需要更加复杂和持久的攻击策略。

\section{单体模型评估的局限性讨论}
\label{sec:limitations}

本章针对医疗大语言模型的单体对抗鲁棒性进行了系统性评估，揭示了模型在数据投毒攻击、后门攻击、提示词注入攻击和越狱攻击等多种威胁下的脆弱性特征。然而，需要指出的是，单体模型的评估框架存在固有的局限性，这些局限性在一定程度上限制了研究结论的普适性和完整性。首先，单体评估框架假设模型在孤立环境中运作，这与实际医疗场景中多模型协作、多系统集成的工作模式存在显著差异。在实际临床工作流中，医疗大语言模型往往作为复杂系统的一部分发挥作用，模型之间、模块之间存在大量的信息交互和依赖关系，这种互联性可能引入单体评估无法覆盖的新型风险。

其次，通用基准数据集虽然能够提供标准化的评估环境，但其与真实临床数据的分布差异可能导致评估结论的偏差。真实世界的电子病历往往包含非规范表述、专业术语缩写、噪声数据和缺失信息，这些特征可能影响模型的防御能力表现。例如，本研究后续章节表明，真实类风湿关节炎数据集上的ASR与MedQuAD数据集存在显著差异，这进一步印证了数据分布对评估结果的影响。此外，现有评估指标在检测特定类型的医疗谬误方面存在盲区。例如，PPL和COH等文本质量指标无法有效检测语义正确但临床错误的输出，这类``高置信度幻觉''在医疗场景中尤为危险。

最后，自动化评估系统虽然能够提高评估效率，但仍然存在``AI评测AI''的循环论证问题。评估模型本身可能存在偏见或安全漏洞，这些缺陷可能传递给评估结果，从而影响结论的可靠性。因此，在未来的研究中，需要引入更多的人工专家复核环节，建立人机协作的评估机制，以确保评估结果的客观性和可信度。

本章实验的局限性还包括以下几个方面：

(1) \textbf{评估偏差来源}：自动化评估系统依赖预定义的关键词匹配和语义相似度计算，可能遗漏某些隐蔽的医疗谬误，或将合理的医学讨论误判为有害输出。

(2) \textbf{结果推广边界}：本研究的实验结果基于英文数据集和特定版本的开源模型，向中文环境和商业模型（如GPT-4、Claude等）的推广需要额外验证。

(3) \textbf{方法论局限}：本章的评估依赖于特定的攻击工具（如AutoPoison、AutoDAN），不同工具可能产生不同的攻击效果，因此结论的普适性受限于所选工具的代表性。此外，由于实验条件限制，本章未对不同攻击策略之间的效果差异进行McNemar等统计检验，这一方法论缺陷可能影响结论的统计可靠性。

(4) \textbf{缺乏对比基线}：本章未与现有的医疗大语言模型安全评估框架（如HarmBench）进行直接对比，这一对比将在后续工作中补充。

(5) \textbf{模型代表性选择的局限性}：在评估攻击隐蔽性阶段（如MT-Bench综合能力测试），受限于API调用成本以及验证攻击隐蔽性特征的实验初衷，本章仅选取了对数据层面攻击最敏感的AlpaCare模型作为测试代表，未将通用能力衰减测试扩展至所有的医疗专用模型及通用基线模型。未来研究可进一步扩展测试的模型范围，以构建更为完备的安全与能力权衡（Safety-Utility Trade-off）对比基线。

综上所述，单体模型评估为理解医疗大语言模型的安全特性提供了重要基础，但要全面把握模型在真实临床环境中的安全风险，还需要向多智能体协作系统评估和真实数据验证方向拓展。

\section{本章小结}
\label{sec:chapter4_summary}

本章基于MLSE框架对医疗大语言模型的单体对抗鲁棒性进行了系统性评估，从数据层攻击和推理层攻击两个维度全面揭示了模型面临的安全威胁。在数据投毒攻击方面，实验结果表明投毒比例与攻击成功率之间存在明确的正相关关系，AlpaCare模型在10\%投毒比例下的ASR达到12.86\%，而BioGPT和BioLlama3模型表现出相对较强的抵御能力。后门攻击实验揭示了特定医学词汇（如``treatment''、``drug''、``method''）作为触发词时能够有效激活植入的后门机制，且攻击的隐蔽性极强，模型的PPL和COH指标在攻击后仍保持在正常范围内。补充实验进一步发现，Oracle模型的能力直接影响投毒攻击的效果，GPT-3.5生成的投毒样本明显优于Gemma-7B，通过角色提示工程可以进一步增强攻击效果。

在提示词操纵攻击方面，医疗专用模型对多种越狱策略的脆弱性普遍高于通用Llama2模型。AlpaCare模型在HarmfulQA数据集上对Combine攻击的ASR高达100\%，这一令人担忧的结果凸显了医疗大语言模型安全对齐机制的不足。AutoDAN越狱攻击的跨样本泛化分析表明，AlpaCare模型存在系统性安全漏洞，泛化ASR在HarmfulQA上达到0.7540。重复攻击尝试次数分析显示，AlpaCare模型在首次尝试时即被攻破，而Asclepius模型需要平均近30次尝试才能成功攻破，表明不同模型的防御能力存在显著差异。

本章的实验结果虽然基于通用基准数据集，但已经充分证明了医疗大语言模型在对抗性攻击下的脆弱性。更为关键的是，传统文本质量指标无法有效检测由攻击引入的医疗谬误，这为后续研究提出了新的方向。本章为第五章基于真实临床数据的多智能体协作系统安全评估奠定了基础，通过建立单体模型的脆弱性基线，为理解更复杂的级联错误传播现象提供了参照框架。
